% page 285 du fac-simile (image p-284 du scan)
% bloc 165.1 x 229.0 mm ; marge gauche 36.9 mm
\pgc{15.240mm}{0.000mm}{
\l{\cel{51}277.}
\l{}
\l{kirsho.\cel{2}Brandio de cerizi, obtenita per fermentacigar cerizi}
\l{\cel{3}nigra ed olia kerni. - DEFIRS.}
\l{}
\l{kirurgio.\cel{2}Parto di la medicin-arto qua traktas precipue la}
\l{\cel{3}lezuri extera (vunduri, osto-rupturi, e c.) e la operaci}
\l{\cel{3}quin igas necesan ula morbi interna. - DEFIRS.}
\l{}
\l{kisar.\cel{2}(\sou{trans.})\cel{2}Presar sua labii, manifeste di afeciono,}
\l{\cel{3}di amo, di amoro, di respekto, adsur la vizajo, la manuo di}
\l{\cel{3}persono. - DE.}
\l{}
\l{\textquotedbl{}\sou{kismet\textquotedbl{}}.\cel{2}(\sou{che la Mohamedisti)}.\cel{2}Fato.}
\l{}
\l{\cel{1}m ro. (\sou{patol}.)\cel{2}Quaza posho sen aperturo, qua kreskas en la}
\l{\cel{3}tisui pro la dilateso di la kondukto-tubi extera di la glan-}
\l{\cel{3}di diversa di qui la orifico divenis obstruktita, e qua kon-}
\l{\cel{3}tenas humori ed altra materii morbala. - EFIRS.}
\l{}
\l{kizelguro.\cel{2}Materio qua konsistas precipue ye silico (til 90\%),}
\l{\cel{3}generale verda, qua devenas de la deskompozeso di diatomei}
\l{\cel{3}qui kaptabis silico ; uzata por facar dinamito, smalio,}
\l{\cel{3}cementi, e c. - IRS.}
\l{}
\l{klado.\cel{2}Skriburo destinita a transskribesor, e quan onu emendas,}
\l{\cel{3}surstrekizas, e c. ante la kopio neta e definitiva. - D.}
\l{}
\l{\sou{klakar}. (\sou{netrans}.)Bruisar pro la intershokeso (instantala)}
\l{\cel{3}di du surfaci plu o min plata. - EF.}
\l{}
\l{klamar.\cel{2}(\sou{trans.)}\cel{3}Dicar (ulo) per vorti pronuncata forte,}
\l{\cel{3}rezonante. - EFI.}
\l{}
\l{klamido. (\sou{epoki antiqua}). Mantelo Greka, Romana, di qua un ek}
\l{\cel{3}la anguli elevesas til la shultro dextra, ube lu an-ligesas}
\l{\cel{3}per agrafo. - DEFIRS.}
\l{}
\l{klano. Quaza tribuo, che la populi Kelta. - DEFIRS.}
\l{}
\l{\sou{klapo}. I. (\sou{tekn.}) Quaza valvo qua apertesas e klozesas quale}
\l{\cel{3}charniro-plako. - II. Peco quan onu povas quaze faldar sur}
\l{\cel{3}altra peco (fixa). - DF.}
\l{}
\l{klara.\cel{2}I. (a) Qua donas o recevas lumo-radii quin nulo inter-}
\l{\cel{3}ceptas ; (b) Di qua nulo desbriligas la pureso. - II. (\sou{metaf.})}
\l{\cel{3}(a) Dicesas pri son-uno, fono, qua ne esas matida, qua reso-}
\l{\cel{3}nas bone en la orelo ; (b) Dicesas pri to quon perceptas}
\l{\cel{3}la spirito - quo ne esas obskura, - quon la spirito komprenas}
\l{\cel{3}bone. - DEFIRS.}
\l{}
\l{klarineto.\cel{2}(\sou{muziko)}. Sufl-instrumento,kun beko e langeto. -}
\l{\cel{3}DEFIS.}
}
